\subsection{Cyclic induction phase multiplicities exclude $R8\_1\_9$}

\begin{lemma}\label{lem:cyclic-induction-phase-mass}
Let $a$ be an invertible object of order $n$ in a pseudo-unitary
fusion category of dimension $D$. Every phase occurring in $I(a)$
is an $n$th root of $\nu_n(a)$, and the total dimension in each of
these $n$ phases, counted with its induction multiplicity, is $D/n$.
\end{lemma}
\begin{proof}
The rotation spaces $\operatorname{Hom}(1,a^m)$ vanish for
$0<m<n$, so $\nu_m(a)=0$. For $m=n$ that space is
one-dimensional, and $\eta=\nu_n(a)$ is a root of unity.
The positive weighted trace
$\operatorname{Tr}(\theta^n I(a))=D\eta$ has maximal absolute
value $D=\dim I(a)$, so every constituent has $\theta^n=\eta$.
Choose $\lambda^n=\eta$ and write $W_j$ for the weight of phase
$\lambda\zeta_n^j$. The zeroth and first $n-1$ trace identities
are $\sum_jW_j=D$ and
$\sum_jW_j\lambda^m\zeta_n^{jm}=0$ for $1\leq m<n$.
Discrete Fourier inversion gives $W_j=D/n$.
\end{proof}

\begin{theorem}\label{thm:r9-center-trace}
The ring $R8\_1\_9$ has no complex fusion categorification.
\end{theorem}
\begin{proof}
By \cite[Lemma~3.1]{SchNear}, take a pseudo-unitary Galois conjugate
and put $d=2+\sqrt5$, $D=8+16d$. The odd-cube pivotal identities
force the second indicators, in source order, to be
$(1,0,0,0,0,1,1,1)$. The pointed generator $a$ has source label1,
with $a^{-1}$ of label2 and $a^2$ of label7. Its associator class
$k\in\mathbb Z/4$ satisfies $\nu_2(a^2)=(-1)^k$, so $k=0$ or2.
The canonical pointed indicators give
$\nu_4(a)=\nu_4(a^{-1})=(-1)^{k/2}$.

The pointed induction Gram block is $8I_4$. Hence a central row
$v$ occurring over $a$ has no other pointed coordinate, and its
dimension has the form $v_a+Bd$, with $B\geq0$.
Lemma~\ref{lem:cyclic-induction-phase-mass} gives weight
$D/4=2+4d$ in each actual phase. Comparing constant coefficients
therefore gives $\sum_{Z\text{ in one phase}}v_a(Z)^2=2$.
Thus $v_1,v_2\leq1$ for every central row. This argument is uniform
in the pointed associator class and does not average away the two
imaginary phases.

The complete formal codegrees are $D,D',8,8,8,8,8,8$;
the six other characters lie in $\mathbb Q(i)$. Actual
character-labelled Galois symmetry transports a forgetful row by
\[
 T(v)=v+(0,0,0,1,1,1,1,0)
 \left(v_0+v_1+v_2+v_7-\frac{v_3+v_4+v_5+v_6}{2}\right),
\]
up to the sign making it nonnegative. Partner rows have equal
normalized relative moments, because their twists are related by
$\sigma^2$ and $\sigma^2$ fixes $\mathbb Q(\sqrt5)$.
The two forced $I(1)$ rows and the six remaining
character-labelled rows leave exactly eight possible low frames.

For arbitrary roots of unity, the first two relative traces lie
in a fixed $17$-vertex hull. It consists of the nine pairs from
orders $1,2,4,5,10,20$, together with the needed negative-half
endpoints: an extra odd factor coprime to $10$ scales both moments
by $\mu(m)/\varphi(m)\in[-1/2,1]$. No root-order bound or
prime-support assumption is used. The full elementary Gauss-sum
proof is supplied in the evidence packet.

For $k=0$, constituents over labels1 and2 have $\theta^4=1$;
for $k=2$ they have $\theta^4=-1$ and their first two relative
traces vanish. Those over label7 have $\theta^2=1$.
Use these phase restrictions, the exact dimension-integrality
condition $D/\dim(Z)$ integral, Galois pairing, and the residual
Gram bound for the entire row orbit. With the Fourier bound,
class zero leaves $64$ row orbits and $405$ moment columns per
frame. The original, weaker bounds suffice for class two and leave
$92$ orbits and $389$ columns.

The Gram and two induction trace identities have $56$ integer
coordinates. The two covectors in
\texttt{r9\_center\_trace\_certificate.json} have nonnegative
pairing with every column, but target pairings $-64$ in class zero
and $-6432$ in class two, for all eight low frames. Their standalone
standard-library verifier reconstructs every bounded domain and
every positive semidefinite condition with exact rational
arithmetic. Neither target is a nonnegative combination of its
columns, excluding both possible pointed classes.
The initial Galois reduction proves the general complex assertion.
\end{proof}
